\newpage
\section{DESK RESEARCH}
Segundo IBGE (2022), o número de adultos com 65 anos ou mais chegou perto de 22,2 milhões, o que representa cerca de 10,9\% da população e um crescimento de 57,4\% em relação a 2010. Também de acordo com o IBGE, a porcentagem de pessoas 60+ que utilizam a internet em seu dia a dia foi de 66\% em 2023.
O observatório Febraban de 2022 realizou uma pesquisa sobre inclusão digital de idosos no país, os resultados apontam um crescimento desta inclusão, 64\% dos cerca de 450 idosos respondentes dizem que fazem uso regular de internet. O celular foi eleito como dispositivo mais utilizado (90\%), e os principais usos são para comunicação (redes sociais e videochamadas), entretenimento (redes digitais e plataformas de streaming), pesquisa de produtos e preços e uso de serviços bancários, principalmente o PIX. As principais dificuldades relatadas é a falta de familiaridade, também relacionada com interfaces complexas e/ou pouco amigáveis, medo de cometer erros durante o uso e insegurança com fraudes e golpes, 70\% responderam que não se sentem seguros online (FEBRABAN, 2022).


No Brasil existem várias iniciativas que promovem a interação social para idosos, com o objetivo de melhorar a qualidade de vida. A principal delas é o SESC, que é referência nacional e está presente em todos os estados. Vale também mencionar outras organizações, como o Instituto Velho Amigo em São Paulo, como foco em idosos em situação de vulnerabilidade social, a Associação dos Amigos dos Idosos do Brasil em Aracaju e o Núcleo de Convivência de Idoso, também em São Paulo


 No Brasil existe um movimento de preocupação com o bem estar das pessoas com mais de 65 anos, podemos ver a partir da quantidade de iniciativas nesse sentido, entretanto a grande maioria destes  projetos são realizados de forma presencial e possuem alcance e capacidade limitados. Agora pensando no que foi dito sobre inclusão digital, idosos cada vez mais acessando e utilizando recursos de internet, existe um potencial de juntar estas soluções. Realizar iniciativas de interação social em ambientes digitais vai aumentar seu alcance e capacidade de atendimento, impactando positivamente muitas mais vidas.
